% answers_short.tex —— 2011 年·山东卷·文科 · 简短答案速查
% 由 tools/make_short_answers.py 自 answers.tex 的 \gksolhead/\gkanswer 提取，
% 并与 2011/2011_数学真题_山东文_答案要点.md 逐题比对后生成，勿手改（改源 MD 后重跑）。

\begin{gknotice}
\textbf{使用说明}：本节只给最终答案便于快速核对，完整推导见解析版。答案为编写者独立推导，\textbf{非官方标准答案}。
\end{gknotice}

\gkpart{选择题}
\begin{center}\begin{tabular}{ccccccccccccc}
\toprule
题号 & 1 & 2 & 3 & 4 & 5 & 6 & 7 & 8 & 9 & 10 & 11 & 12 \\
\midrule
答案 & A & D & D & C & A & B & B & B & C & C & A & D \\
\bottomrule
\end{tabular}\end{center}

\gkpart{填空题}
\begin{itemize}[leftmargin=2.6em]
  \item[\textbf{13.}] $16$。
  \item[\textbf{14.}] $68$。
  \item[\textbf{15.}] $\dfrac{x^2}{4}-\dfrac{y^2}{3}=1$。
  \item[\textbf{16.}] $2$。
\end{itemize}

\gkpart{解答题}
\begin{itemize}[leftmargin=2.6em]
  \item[\textbf{17.}] (1) $2$；(2) $b=2$
  \item[\textbf{18.}] (1) $\dfrac49$；(2) $\dfrac25$
  \item[\textbf{19.}] (1) 证明见解析；(2) 证明见解析
  \item[\textbf{20.}] (1) $a_n=2\cdot3^{n-1}$；(2) $S_{2n}=3^{2n}-1+n\ln3$
  \item[\textbf{21.}] (1) $y=\dfrac{160\pi}{r}+4\pi(c-2)r^2$，定义域 $0<r\leqslant2$；(2) 当 $3<c\leqslant4.5$ 时 $r=2$；当 $c>4.5$ 时 $r=\sqrt[3]{\dfrac{20}{c-2}}$
  \item[\textbf{22.}] (1) $2$；(2) ① 直线 $l$ 过定点 $(-1,0)$；② 点 $B$，$G$ 能关于 $x$ 轴对称，此时 $\triangle ABG$ 的外接圆方程为 $\left(x+\dfrac12\right)^2+y^2=\dfrac54$
\end{itemize}
